% \section{Technical appendices and supplementary material}
% Technical appendices with additional results, figures, graphs, and proofs may be submitted with the paper submission before the full submission deadline (see above). You can upload a ZIP file for videos or code, but do not upload a separate PDF file for the appendix. There is no page limit for the technical appendices. 

% Note: Think of the appendix as ``optional reading'' for reviewers. The paper must be able to stand alone without the appendix; for example, adding critical experiments that support the main claims to an appendix is inappropriate. 


% ============================================================
% Supplementary Appendix for EFFORT
% ============================================================
% All values below were verified against the training log, phase summary, and
% code before release. Sources map to the release package as:
%   segmentation  -> code/phase_detector.py + data/phase_log.csv
%   sonification  -> code/gen_sonification.py
%   effort words  -> the rule in this appendix + the phase_log recomputation
%   alignment     -> code/build_stage_timeline.py (analyze_audio)
%   raw log check -> data/training_log.csv.gz
%   hyperparams   -> code/config/lodge_coarse_finedance.yaml
% A few items are genuinely external (the Suno model version and per-generation
% id, which the service does not expose). Those are stated honestly rather than
% invented. There are no remaining TODOs.
%
% Expected bibliography keys:
%   li2024lodge, li2023finedance, laban1947effort, suno2026
%
% Recommended package in the main document:
%   \usepackage{amsmath}
% ============================================================

% \appendix

\section{Model training and dataset}

\emph{Effort} uses Lodge \citep{li2024lodge}, a two-stage coarse-to-fine diffusion network that generates dance motion conditioned on music, trained from scratch for this project.
The trace analyzed here comes from the Global module, the coarse stage of the network.
Music enters the model as 35-dimensional handcrafted audio features, and motion is represented as a 139-dimensional feature vector over 22 joints.
The trace therefore records the optimization history of this specific run.
The training data is FineDance \citep{li2023finedance}, a motion-capture dataset of paired music and full-body 3D dance recordings performed by twenty-seven professional dancers, used in its standard training split.
% In the version of the dataset available to us, the performers are credited collectively, without individual names attached to the motion sequences.
Training ran for 3,500 epochs on a single NVIDIA H100, about 51 hours of active GPU time.
The released configuration records the optimizer settings: Adan with a learning rate of \(4\times10^{-4}\) held constant, batch size 128 with gradient accumulation of three for an effective batch of 384.
We disabled gradient clipping, so the logged gradient norm is the raw norm.
A logging callback writes one row per batch, giving 155,232 rows over 52,516 optimizer steps.
Each row records the step, epoch, total and component losses, gradient-norm statistics, the learning rate, and binary indicators for loss spikes, plateaus, and exploding or vanishing gradients, and every indicator is false in every row.

The artwork uses three signals from the log: training loss, mean gradient norm, and gradient-norm volatility.
Volatility is measured two ways.
At the phase level, it is the standard deviation of the gradient norm over the phase's downsampled rows, and the ratio of that standard deviation to the phase mean drives the Effort mapping in Section~\ref{app:effort}.
At the per-batch level, the log stores a sign-change rate, the fraction of gradient coordinates whose sign flips between consecutive batches, and the sonifier in Section~\ref{app:music} uses it for rhythm.

A change-point detector divides the run into fourteen contiguous phases.
The per-batch log is downsampled to every fifteenth row in acquisition order, and the PELT algorithm from the \texttt{ruptures} library is run on the smoothed loss and on the gradient norm with a radial-basis-function cost, a penalty of 10, a merge window of 200 rows, and a minimum segment length of 500 rows, all in downsampled rows.
The two sets of breakpoints are merged and the minimum-length constraint is enforced, which yields fourteen phases.
Table~\ref{tab:phases} lists the per-phase statistics.
The run opens with a phase of large, unstable gradients, followed by thirteen lower-amplitude phases of gradual convergence in which the mean loss falls from about 0.116 to about 0.0098.
The artistic transformations described below operate on this continuous training history.

\section{Translating the training trace into music}
\label{app:music}

The first translation turns the per-batch training log into a musical score through an artist-designed, rule-based procedure.
Given the same log and settings, the procedure regenerates the same symbolic score.
The score is written for a mechanical player piano with a small chamber ensemble, in D natural minor, in 4/4, at 70 beats per minute.
The player piano carries the melody, strings and a low synthesizer drone support it, and celesta, harp, contrabass, and string pad enter occasionally.
Four quantities from the log control four musical parameters, each normalized against this run's own range.
\begin{enumerate}
\item \textbf{Register.} The gradient norm sets the register of the melody and accompaniment. Larger gradients move the material upward.
\item \textbf{Rhythmic density.} The gradient sign-change rate sets the number and spacing of note onsets. Stable intervals give sparse, sustained notes, and volatile intervals give denser subdivisions and syncopation.
\item \textbf{Tremolo and chromatic inflection.} Bars with a change in loss above the 82nd percentile across the piece receive a string tremolo and unresolved chromatic neighbor tones. The lines are otherwise diatonic.
\item \textbf{Spectral brightness.} The gradient level sets the brightness of the synthesizer pad, so the pad darkens after the high-gradient opening.
\end{enumerate}

Data controls register, rhythmic density, the accents, and the phase durations.
The key, the harmonic progressions, the instrumentation, and the five-part form are artist-designed.

The fourteen training phases are laid out in training order on one linear clock.
Each phase receives bars in proportion to its length, with a floor of two bars to keep every phase audible.
Music time therefore follows training time at bar-level resolution.
Over this clock the instrumentation follows a fixed five-part shape, a piano-led opening, a build, a short caesura, a denser passage, and a decay as the layers drop out.
The score ends on an open fifth, D and A, and its final four bars slow from 70 to about 38 beats per minute.
The Effort words in Section~\ref{app:effort} are indexed by the same fourteen phases, so the two translations share one clock.

\section{Translating the training trace into Laban Effort words}
\label{app:effort}

The second translation gives one Laban Effort word to each of the fourteen phases through the three Effort factors of \textit{Weight}, \textit{Time}, and \textit{Space} \citep{laban1947effort}.
The mapping is an artist-designed convention applied to the statistics of this run.

For phase \(p\), let \(\mu_p\) be its mean gradient norm, \(\sigma_p\) the standard deviation of the gradient norm, and \(\tau_p\) its temporal trend, the slope of a least-squares fit of the gradient norm over the phase.
Let \(m_\mu\) and \(m_\tau\) be the medians of \(\mu_p\) and \(|\tau_p|\) over the fourteen phases.
\textit{Weight} is strong when \(\mu_p \geq m_\mu\) and light otherwise.
\textit{Time} is sudden when \(\sigma_p \geq \mu_p\) and sustained otherwise.
\textit{Space} is direct when \(|\tau_p| \geq m_\tau\) and indirect otherwise.
\textit{Weight} and \textit{Space} compare a phase with the whole run, and \textit{Time} compares a phase's variability with its own mean.
The three binary factors give the eight Effort words in Table~\ref{tab:effort}.

\begin{table}[h]
  \centering
  \begin{tabular}{llll}
    \hline
    Weight & Time & Space & Effort \\
    \hline
    light  & sustained & indirect & \textit{Float} \\
    light  & sustained & direct   & \textit{Glide} \\
    light  & sudden    & direct   & \textit{Dab} \\
    light  & sudden    & indirect & \textit{Flick} \\
    strong & sustained & direct   & \textit{Press} \\
    strong & sustained & indirect & \textit{Wring} \\
    strong & sudden    & direct   & \textit{Punch} \\
    strong & sudden    & indirect & \textit{Slash} \\
    \hline
  \end{tabular}
  \caption{Artist-designed mapping from the three binary Laban Effort factors
  to eight Effort words. Each training phase is assigned one combination.}
  \label{tab:effort}
\end{table}

The first phase, with large and unstable gradients and a steep downward trend, is strong, sudden, and direct, the Effort \textit{Punch}.
It is the only sudden phase, since every later phase has a standard deviation below its mean.
The run then passes through \textit{Press}, settles into \textit{Float} and \textit{Glide} through its low-gradient middle, and ends in \textit{Press} and \textit{Wring} as the gradient rises again.
Table~\ref{tab:phases} lists the full assignment.
The resulting Effort sequence is a choreographic interpretation constructed through our selection of signals, segmentation, thresholds, and factor mapping.
In the performance the dancer receives each word as an open score. The word sets the quality of a movement, and the pose and trajectory remain their choice.

\begin{table*}[t]
  \centering
  \small
  \begin{tabular}{rrrrrrllll}
    \hline
    Phase & Start step & Epoch & $\mu_p$ & $\sigma_p$ & $\tau_p$ & Weight & Time & Space & Effort \\
    \hline
    1  & 0     & 0    & 2.741 & 5.017 & $-1.6\times10^{-2}$ & strong & sudden    & direct   & \textit{Punch} \\
    2  & 2593  & 172  & 0.343 & 0.273 & $-2.3\times10^{-4}$ & strong & sustained & direct   & \textit{Press} \\
    3  & 5610  & 374  & 0.235 & 0.175 & $+9.4\times10^{-6}$ & light  & sustained & indirect & \textit{Float} \\
    4  & 8704  & 580  & 0.233 & 0.200 & $-9.8\times10^{-5}$ & light  & sustained & direct   & \textit{Glide} \\
    5  & 11260 & 750  & 0.248 & 0.198 & $+1.0\times10^{-4}$ & light  & sustained & direct   & \textit{Glide} \\
    6  & 15019 & 1001 & 0.272 & 0.170 & $+3.2\times10^{-5}$ & light  & sustained & indirect & \textit{Float} \\
    7  & 18292 & 1219 & 0.283 & 0.180 & $-1.4\times10^{-5}$ & light  & sustained & indirect & \textit{Float} \\
    8  & 21155 & 1410 & 0.313 & 0.245 & $+1.3\times10^{-4}$ & light  & sustained & direct   & \textit{Glide} \\
    9  & 25016 & 1667 & 0.341 & 0.204 & $+5.3\times10^{-5}$ & light  & sustained & direct   & \textit{Glide} \\
    10 & 28902 & 1926 & 0.370 & 0.218 & $+6.5\times10^{-5}$ & strong & sustained & direct   & \textit{Press} \\
    11 & 33530 & 2235 & 0.386 & 0.221 & $+4.6\times10^{-5}$ & strong & sustained & indirect & \textit{Wring} \\
    12 & 36394 & 2426 & 0.398 & 0.247 & $+2.2\times10^{-5}$ & strong & sustained & indirect & \textit{Wring} \\
    13 & 43374 & 2891 & 0.379 & 0.209 & $+1.3\times10^{-5}$ & strong & sustained & indirect & \textit{Wring} \\
    14 & 45931 & 3062 & 0.370 & 0.207 & $+3.2\times10^{-6}$ & strong & sustained & indirect & \textit{Wring} \\
    \hline
  \end{tabular}
  \caption{The fourteen phases of the training run. Start step is the first
  logged optimizer step of each phase. Epoch is the approximate epoch at
  that boundary. The phases are contiguous, so a phase runs to the next phase's
  start. The remaining columns give the mean gradient norm $\mu_p$, the
  gradient-norm standard deviation $\sigma_p$, the temporal trend $\tau_p$, the
  three binary factors, and the resulting Effort word. The median of $\mu_p$ is $0.342$ and the median
  of $|\tau_p|$ is $4.9\times10^{-5}$.}
  \label{tab:phases}
\end{table*}

\section{Suno re-performance and temporal alignment}

The rule-based score fixes the melody, harmony, form, phase durations, and tempo.
For the video we passed its rendering to Suno \citep{suno2026}, a generative music service, in its Cover mode.
This step is a generative re-performance outside the deterministic mapping, and Suno can return a different cover on each request, so we keep the selected outputs as fixed artifacts.

The presented cover was generated with model V5.5 as an instrumental, with Audio Influence 85\%, Style Influence 15\%, and Weirdness 5\%, from the following prompt.
We generated six candidates with the same prompt and settings and selected the take closest to the source player-piano motif and its restrained dynamic.
The file is used as Suno delivered it.

\begin{quote}
\small
COVER FAITHFULLY --- preserve the uploaded melody, chord changes, tempo
(keep the final ritardando), and section structure. Keep the dynamics
restrained from start to end: quiet, chamber-scale, mezzo-piano at most.
The ``climax'' is only a slight thickening of texture, then it thins again.
A small, cold ensemble: a mechanical player piano stating the melody with
precise, unhurried articulation; one solo cello answering low and dry; a
faint, even eighth-note pulse in muted low strings, felt more than heard;
sparse high violin harmonics for air. Wide space between phrases, generous
silence, cool and exact intonation. D minor, approximately 70 BPM, 4/4.
The ending slows and settles on a bare open fifth, unresolved, fading into
a dry, cold room tail. Mood: wintry, elegant, austere, lucid, still;
mechanical calm rather than passion. Instrumental only, close-miked modern
classical recording, high detail, quiet mastering.
\end{quote}

The cover changes the timing of the score, so we align the visual materials to it with chroma-based dynamic time warping, a monotone map from score time to cover time.
The plotted curves, bar and phase boundaries, and Effort-word changes all follow the timing of the audio in the video.

The dance is conditioned on a second cover of the same score, arranged in the style of Kunqu opera accompaniment.
The model produces more coherent motion from this cover than from the sparse texture of the presented cover, and its movement carries the stillness and austerity closest to the finished work.
This cover keeps the key, phrase structure, and five-part form of the score at close to the same tempo, and differs from the presented cover in instrumentation and surface character.
The generated dance is aligned to the presented cover with the same warping.
It was generated with model V4.5, Audio Influence 85\%, and Weirdness 0\%, from the following prompt.

\begin{quote}
\small
COVER, REARRANGED IN THE STYLE OF KUNQU OPERA ACCOMPANIMENT. Keep the
source's slow tempo, minor-key character, and melodic contour. Replace the
Western orchestration with traditional Kunqu instrumentation. D minor with
pentatonic flavor, approximately 65 BPM, 4/4. Production: bamboo flute (qudi
or dizi) carrying the melodic line, sheng (mouth organ) sustained drones,
pipa or sanxian plucked accents on phrase boundaries, soft guzheng arpeggios,
light wooden percussion (wooden clappers and paigu on downbeats), small bells.
Subtle wordless female vocalise in Kunqu falsetto style --- breathy, soft,
ornamented --- hovering as background, sparse, only in climax and outro.
Natural temple-hall reverb. Mood: meditative, classical Chinese, sacred
lament, slow and flowing. Suitable as score for Kunqu and classical Chinese
stage choreography. Mostly instrumental.
\end{quote}

\section{Reproducibility and retained artifacts}

The release at \url{https://github.com/ahugh19/Effort} contains the per-batch training log, the phase summary, the symbolic score and its mastered render, the two Suno covers, the alignment timeline, and the code for each deterministic step.
The pipeline runs in the following order. Steps 1 and 5 are outside the release, and their outputs are kept as fixed files.
\begin{enumerate}
\item Train Lodge on FineDance and write the per-batch log.
\item Segment the log into fourteen phases.
\item Generate the symbolic score and, optionally, a dry preview render.
\item Assign one Effort word to each phase.
\item Generate and select the Suno cover.
\item Align the score to the presented cover by dynamic time warping.
\item Render the curve, word, dancer, and skeleton layers on the aligned timeline.
\end{enumerate}
The exact parameters of steps 2, 3, 4, and 6 are recorded in the released scripts.
The mastered render stays a fixed file because mastering happened outside the package, and the covers stay fixed files because the service returns a different result on each request.
The FineDance recordings remain with their authors \citep{li2023finedance} under the dataset license, and the released material is provided under the license stated in the release README.